% sections/sec2_bilinear.tex
\section{双线性映射}
\label{sec:bilinear}

我们在前一节中已经将 $M \otimes_R N$ 的元素描述为服从规则~\eqref{eq:1.2} 与~\eqref{eq:1.3} 的和式~\eqref{eq:1.4}.其核心意图在于使得 $M \otimes_R N$ 成为满足规则~\eqref{eq:1.2} 与~\eqref{eq:1.3} 的``最自由''对象.而规则~\eqref{eq:1.2} 与~\eqref{eq:1.3} 的本质就是双线性（bilinearity）.这究竟是什么含义？

设 $M$、$N$ 与 $P$ 为 $R$-模.如果函数 $B \colon M \times N \to P$ 在固定另一个变元时对每个变元都是线性的（即 $R$-线性的），则称 $B$ 为\newterm{bilinear-map}{双线性映射}{bilinear map}（或双线性函数）：
\begin{align}
  B(m_1 + m_2, n) &= B(m_1, n) + B(m_2, n), & B(rm, n) &= rB(m, n), \label{eq:2.1} \\
  B(m, n_1 + n_2) &= B(m, n_1) + B(m, n_2), & B(m, rn) &= rB(m, n). \label{eq:2.2}
\end{align}
因此对每个 $n \in N$，$B(-, n)$ 是一个线性映射 $M \to P$；对每个 $m \in M$，$B(m, -)$ 是一个线性映射 $N \to P$.特别地，总有 $B(0, n) = 0$ 以及 $B(m, 0) = 0$.

以下是双线性映射的一些常见实例：
\begin{enumerate}
  \item $\R^n$ 上的点积 $v \cdot w$ 是双线性函数 $\R^n \times \R^n \to \R$.更一般地，对任意矩阵 $A \in \mathrm{M}_n(\R)$，函数 $\langle v, w \rangle = v \cdot Aw$ 都是双线性映射 $\R^n \times \R^n \to \R$.
  \item 矩阵乘法 $\mathrm{M}_{m,n}(\R) \times \mathrm{M}_{n,p}(\R) \to \mathrm{M}_{m,p}(\R)$ 是双线性的.点积是 $m = p = 1$ 的特例（此时将 $v \cdot w$ 写作 $v^\top w$）.
  \item 叉积 $v \times w$ 是双线性函数 $\R^3 \times \R^3 \to \R^3$.
  \item 行列式 $\det \colon \mathrm{M}_2(R) \to R$ 作为矩阵两列的函数是双线性的.
  \item 对于 $R$-模 $M$，标量乘法 $R \times M \to M$ 是双线性的.
  \item 环自身的乘法 $R \times R \to R$ 是双线性的.
  \item 设 $M$ 的对偶模为 $M^\vee = \Hom_R(M, R)$.由 $(\varphi, m) \mapsto \varphi(m)$ 给出的对偶配对 $M^\vee \times M \to R$ 是双线性的.
  \item 对于 $\varphi \in M^\vee$ 与 $\psi \in N^\vee$，映射 $M \times N \to R$（其中 $(m, n) \mapsto \varphi(m)\psi(n)$）是双线性的.
  \item 若 $M \times N \xrightarrow{B} P$ 是双线性的，且 $P \xrightarrow{L} Q$ 是线性的，则复合映射 $M \times N \xrightarrow{L \circ B} Q$ 也是双线性的.（这是一个至关重要的性质，请务必验证！）
  \item 根据第~\ref{sec:intro}~节的讨论，表达式 $m \otimes n$ 在 $m$ 与 $n$ 上应当是双线性的.换言之，我们希望由 $(m, n) \mapsto m \otimes n$ 给出的函数 $M \times N \to M \otimes_R N$ 是双线性的.
\end{enumerate}

以下是几个双元函数但\textbf{并非}双线性映射的典型反例：
\begin{enumerate}
  \item 对于 $R$-模 $M$，由 $(m, m') \mapsto m + m'$ 给出的加法运算 $M \times M \to M$ 通常不是双线性的：当 $m'$ 固定时它通常对 $m$ 不是加性的（即一般而言 $(m_1 + m_2) + m' \neq (m_1 + m') + (m_2 + m')$），当 $m$ 固定时它对 $m'$ 亦非加性.
  \item 对于 $\varphi \in M^\vee$ 与 $\psi \in N^\vee$，由 $(m, n) \mapsto \varphi(m) + \psi(n)$ 给出的和式 $M \times N \to R$ 通常不是双线性的.
  \item 将 $\mathrm{M}_n(\C)$ 视为 $\C$-向量空间.由 $(A, B) \mapsto A\bar{B}$ 给出的函数 $\mathrm{M}_n(\C) \times \mathrm{M}_n(\C) \to \mathrm{M}_n(\C)$ 不是双线性的.它是双加性的（即当固定另一个分量时对每个分量均可加），但观察标量乘法在第二个分量上的表现：对 $z \in \C$，$A\overline{zB} = \bar{z}(A\bar{B})$，而不是 $z(A\bar{B})$.
\end{enumerate}

为什么我们要把双线性函数写成 $B \colon M \times N \to P$，而不是写成 $B \colon M \oplus N \to P$？因为我们并没有在定义域上使用模的直和加法结构.这就是为什么对于加法 $R \times R \to R$ 与乘法 $R \times R \to R$，我们可以把加法写成 $R \oplus R \to R$（它在直和 $R \oplus R$ 上是线性的），但绝不能对乘法这样做：乘法是双线性的，却不是线性的（例如一般而言 $(r + r')(s + s') \neq rs + r's'$）.线性函数是广义的加法，而双线性函数是广义的乘法.切勿将 $M \times N$ 上的双线性函数与 $M \oplus N$ 上的线性函数混淆.

对于 $R$-模 $M_1, \dots, M_k$，如果函数 $f \colon M_1 \times \dots \times M_k \to M$ 在固定其他坐标时对每个 $m_i$ 都是线性的（即 $R$-线性的），则称 $f$ 为\newterm{multilinear-map}{多重线性映射}{multilinear map}（或 $k$-重线性映射）.因此 $2$-重线性即指双线性.以下是多重线性函数的一些例子：
\begin{enumerate}
  \item 标量三重积（混合积）$u \cdot (v \times w)$ 是三线性函数 $\R^3 \times \R^3 \times \R^3 \to \R$.
  \item 函数 $f(u, v, w) = (u \cdot v)w$ 是三线性函数 $\R^n \times \R^n \times \R^n \to \R^n$.
  \item 由 $(\varphi, m, n) \mapsto \varphi(m)n$ 给出的函数 $M^\vee \times M \times N \to N$ 是三线性的.
  \item 若 $B \colon M \times N \to P$ 与 $B' \colon P \times Q \to T$ 均为双线性的，则由 $(m, n, q) \mapsto B'(B(m, n), q)$ 给出的映射 $M \times N \times Q \to T$ 是三线性的.
  \item 具有 $k$ 个因子的乘法 $R \times \dots \times R \to R$ 是 $k$-重线性的.
  \item 行列式 $\det \colon \mathrm{M}_n(R) \to R$ 作为矩阵各列的函数是 $n$-重线性的.
  \item 若 $M_1 \times \dots \times M_k \xrightarrow{f} M$ 是 $k$-重线性的，且 $M \xrightarrow{L} N$ 是线性的，则复合映射 $M_1 \times \dots \times M_k \xrightarrow{L \circ f} N$ 也是 $k$-重线性的.
  \item 对于光滑映射 $f \colon \R^m \to \R^n$ 及点 $x \in \R^m$，$k$ 阶导数 $(D^k f)_x \colon (\R^m)^k \to \R^n$ 是 $k$-重线性的.（若不熟悉将一阶导数视作线性映射的观点，可忽略此例.）
\end{enumerate}

从 $M$ 到 $N$ 的所有 $R$-线性映射在函数的点态加法与标量数乘下构成一个 $R$-模 $\Hom_R(M, N)$.类似地，从 $M \times N$ 到 $P$ 的所有 $R$-双线性映射以相同方式构成一个 $R$-模 $\mathrm{Bil}_R(M, N; P)$.然而，与线性映射相比，双线性映射缺失了若干关键特征：
\begin{enumerate}
  \item 双线性映射 $M \times N \to P$ 不存在``核''（kernel）的概念，因为定义域 $M \times N$ 并非一个模（我们没有赋予其模运算使其保持线性）.
  \item 双线性映射 $M \times N \to P$ 的像集不一定构成子模.
\end{enumerate}

\begin{example}\label{ex:rank-one}
定义 $B \colon \R^n \times \R^n \to \mathrm{M}_n(\R)$ 为 $B(v, w) = v w^\top$，其中 $v$ 与 $w$ 均为列向量，因此 $v w^\top$ 是一个 $n \times n$ 矩阵.例如当 $n = 2$ 时，
\[
  B\left(\begin{pmatrix} a_1 \\ a_2 \end{pmatrix}, \begin{pmatrix} b_1 \\ b_2 \end{pmatrix}\right) = \begin{pmatrix} a_1 \\ a_2 \end{pmatrix} (b_1, b_2) = \begin{pmatrix} a_1 b_1 & a_1 b_2 \\ a_2 b_1 & a_2 b_2 \end{pmatrix}.
\]
一般地，若用 $\R^n$ 的标准基表示 $v = \sum a_i e_i$ 与 $w = \sum b_j e_j$，则 $v w^\top$ 就是 $n \times n$ 矩阵 $(a_i b_j)$.该矩阵关于 $v$ 和 $w$ 分别都是 $\R$-线性的，所以 $B$ 是双线性的.

然而当 $n \geq 2$ 时，$B$ 的像在加法下并不封闭，因此它不是 $\mathrm{M}_n(\R)$ 的子空间.原因何在？由于矩阵 $B(v, w)$ 的每一列都是列向量 $v$ 的标量倍数，所以其秩至多为 $1$（非零时秩恰为 $1$，零时秩为 $0$）.考虑矩阵
\[
  B(e_1, e_1) + B(e_2, e_2) = e_1 e_1^\top + e_2 e_2^\top = \begin{pmatrix}
    1 & 0 & \dots & 0 \\
    0 & 1 & \dots & 0 \\
    \vdots & \vdots & \ddots & \vdots \\
    0 & 0 & \dots & 0
  \end{pmatrix},
\]
该矩阵的像空间是 $2$ 维的（秩为 $2$），因此对于 $\R^n$ 中的任何向量 $v$ 与 $w$，都不可能有 $B(e_1, e_1) + B(e_2, e_2) = B(v, w)$.（类似地，$\sum_{i=1}^n B(e_i, e_i)$ 是 $n \times n$ 单位矩阵，它绝不能写成单个矩阵 $B(v, w)$ 的形式.）
\end{example}
